\documentclass{article}
\usepackage[utf8]{inputenc}
\usepackage[T1]{fontenc}
\usepackage{geometry}
\geometry{
    a4paper,
    left=2.5cm,
    right=2.5cm,
    top=2.5cm,
    bottom=2.5cm
}
\usepackage{graphicx}
\usepackage{float}
\usepackage{hyperref}
\usepackage[english]{babel}
\usepackage{amsmath}
\usepackage{amssymb}
\usepackage{array}
\usepackage{booktabs}
\usepackage{caption}
\usepackage{subcaption}
\usepackage{enumitem}
\usepackage{fancyhdr}
\usepackage{titlesec}
\usepackage{url}
\usepackage{textcomp}
\usepackage{xcolor}
\usepackage{listings}
\usepackage{multirow}
\usepackage{longtable}
\usepackage{tabularx}
\usepackage{times}

% Set up page style
\pagestyle{fancy}
\fancyhf{}
\rfoot{\thepage}
\renewcommand{\headrulewidth}{0pt}

% Title formatting
\titleformat{\chapter}[display]{\normalfont\bfseries}{\chaptertitle\thechapter}{0pt}{\Huge}
\titlespacing*{\chapter}{0pt}{50pt}{40pt}

% Document info
\title{Template Document}
\author{Academic Support Agent}
\date{\today}

\begin{document}
\maketitle

%
	%
	%%%%%%%%  The "To" address goes here.
	%
	%Also, you can use this
	% \begin{letter}{
	% 		Target University\\ 
	% 		Some Addresss\\ 
	% 		SomeTown, SomeState 					       				  		 ~~SomeZip
	% 	}
\begin{letter}{To whom it may concern,}

		% This line sets up the return address to the right-side of the OSU logo.  The location is set with absolute node addresses using ``\tikz''.  It can still be a bit fussy, and you may need to alter this a little to get things to look right.  The bit that changes the position are the numbers in parentheses ``at (14.2,2.7)''
		%
		\begin{tikzpicture}[remember picture,overlay,,every node/.style={anchor=center}]
		\node[text width=7cm] at (page cs:0.5,0.73){\small \newaddress};
		\end{tikzpicture} 
		
		%%%%%%  The ``opening'' is just the method of address you would like to use at the start of the letter.
		%
        \opening{ }

		PARAGRAPH ONE: State the reason for the letter, name the position or type of work you are applying for and identify the source from which you learned of the opening (i.e. career development center, newspaper, employment service, personal contact).
		
		PARAGRAPH TWO: Indicate why you are interested in the position, the company, its products, services - above all, stress what you can do for the employer. If you are a recent graduate, explain how your academic background makes you a qualified candidate for the position. If you have practical work experience, point out specific achievements or unique qualifications. Try not to repeat the same information the reader will find in the resume. Refer the reader to the enclosed resume or application which summarizes your qualifications, training, and experiences. The purpose of this section is to strengthen your resume by providing details which bring your experiences to life. 
		
		PARAGRAPH THREE:End the letter by thanking the employer for taking time to consider your credentials. 

		%%%%%%% ``closing'' sets the sign-off line.
		\closing{Kind regards,}
		
		% Comment out/in the lines below as necessary
		%\encl{If an enclosure is provided, let them know what it is.}
		
		%\ps{A postscript if that is a thing you do.}
		
		%\cc{Someone Who Cares (and is copied).}
		
	\end{letter}
\end{document}